% ==============================================================================
% SECCIÓN 3: DEFINICIÓN DE LA VISIÓN DEL SISTEMA (PARTE 2 DE LA GUÍA)
% ==============================================================================

\section{Parte 2: Definición de la Visión del Producto}

\subsection{Plantilla Canónica de Geoffrey Moore}
La visión estratégica de \textbf{SIGC-CUSCO} se formaliza mediante la plantilla de posicionamiento de producto:

\begin{tcolorbox}[colback=blue!4!white, colframe=blue!70!black, title=\textbf{Declaración Oficial de Visión -- SIGC-CUSCO}, arc=1.5mm]
\begin{description}\setlength\itemsep{0.2em}
    \item[\textbf{Para:}] Municipalidades distritales y empresas privadas (turismo, hotelería y comercio del Cusco) que gestionan actividades de capacitación y actualización técnica.
    \item[\textbf{El sistema:}] \textbf{SIGC-CUSCO (Sistema Integral de Gestión de Capacitaciones)}.
    \item[\textbf{Que permite:}] Centralizar en una plataforma web todo el proceso formativo: convocatoria, inscripción digital, control de asistencia en tiempo real, evaluaciones y emisión automática de certificados con código QR verificable.
    \item[\textbf{A diferencia de:}] Hojas de cálculo de Excel aisladas, listas de asistencia en papel propensas a pérdida o fraude, y diseño manual de diplomas sin validación pública.
    \item[\textbf{Nuestro producto:}] Ofrece un sistema multientidad, escalable y seguro que reduce en más del 80\,\% los tiempos administrativos y garantiza autenticidad jurídica inmediata en cada certificado.
\end{description}
\end{tcolorbox}

\subsection{Delimitación del Alcance (Matriz Es / No Es / Hace / No Hace)}

\begin{table}[H]
\centering
\small
\renewcommand{\arraystretch}{1.2}
\begin{tabularx}{\textwidth}{XX}
\toprule
\textbf{El Sistema SIGC-CUSCO SÍ ES / SÍ HACE} & \textbf{El Sistema SIGC-CUSCO NO ES / NO HACE} \\
\midrule
\textbullet~Plataforma web de gestión integral de capacitaciones. & \textbullet~NO es un LMS complejo para carreras de pregrado (como Blackboard). \\
\textbullet~Inscripción pública y corporativa con validación de DNI. & \textbullet~NO procesa pasarelas de pago bancarias directas en el MVP inicial. \\
\textbullet~Control de asistencia digital mediante código QR / PIN. & \textbullet~NO transmite streaming de video en vivo (enlaza a Zoom/Meet). \\
\textbullet~Generación de certificados PDF con QR y hash SHA-256. & \textbullet~NO es una red social de búsqueda de empleo. \\
\textbullet~Reportes y padrones oficiales de participantes y notas. & \textbullet~NO emite títulos a nombre de la Nación sin aval oficial. \\
\bottomrule
\end{tabularx}
\caption{Delimitación del alcance funcional del proyecto SIGC-CUSCO.}
\label{tab:alcance_frontera}
\end{table}

\subsection{Objetivos Estratégicos SMART}
\begin{itemize}\setlength\itemsep{0.2em}
    \item \textbf{Específico (S):} Implementar los módulos de Instituciones, Capacitaciones, Asistencia QR, Evaluación y Certificación Digital.
    \item \textbf{Medible (M):} Reducir el tiempo de emisión de diplomas de 7 días hábiles a menos de 10 segundos tras finalizar el curso.
    \item \textbf{Alcanzable (A):} Desarrollar un MVP funcional multitenant dentro del cronograma del curso de Ingeniería de Software I.
    \item \textbf{Relevante (R):} Optimizar la gestión pública local en Cusco y elevar la competitividad de las empresas turísticas.
    \item \textbf{Temporal (T):} Completar la fase de \textit{Inception} y requerimientos en las primeras semanas del ciclo 2026-II.
\end{itemize}
